%% Data and Preprocessing
\section*{Data and preprocessing}

We sample 100 molecules from ChEMBL with pre-stratified difficulty. We first run PrexSyn on a larger pool to compute baseline quality scores, defined as the ECFP4 Tanimoto similarity between the best candidate and the specification fingerprint, then draw approximately 25 molecules from each of four quality bins ($<$0.5, 0.5--0.7, 0.7--0.85, 0.85--1.0), ensuring balanced representation across difficulty levels.

For each specification, PrexSyn generates 256 candidates. The top-scoring candidate ($K=1$) seeds all four modification methods, maintaining clean causal attribution. Each method generates up to $N$ variants, where $N$ is set by a pilot saturation study (5 specs $\times$ 4 methods $\times$ 1,000 variants) that identifies the diversity plateau beyond which additional variants yield diminishing structural novelty. We deduplicate variants by canonical SMILES before checking synthesizability via AiZynthFinder at search depth 1. Because depth 2 is $\sim$5$\times$ slower, we run a depth-2 sensitivity analysis on a 10\% subsample only. The estimated total load is $\sim$80,000 unique molecules.

We train no models from scratch. CReM, mmpdb, and JT-VAE operate as pre-trained or pre-configured tools. LibINVENT fine-tunes a pretrained prior via reinforcement learning to steer decoration toward each specification's target properties. CReM's fragment database derives from ChEMBL, the same source as our property specifications; ChEMBL molecules serve only as property-profile sources, not reconstruction targets.
